% Folder 74, batch 06 (archivists' pages 101-120).
% Pass: Opus 5.5 (claude-opus-5-5), 2026-10-03 — first pass, unchecked against the pages by a human.
%
% Notation fixed for this batch: tilde objects (\widetilde{E}_a, \widetilde\Gamma)
% for the double coverings, \underline{t} for a triangle taken as an object,
% \mathcal{G} for the automorphism group of the trialitarian graph (S,A),
% \mathcal{G}^{\circ} for the subgroup generated by the symmetries \sigma_{\underline t}.
% Page 112 was read turned through 180 degrees. No page of the batch is skipped.
\documentclass[11pt,a4paper]{article}
\input{../preamble/grothendieck}

\folder{74}
\batch{6}
\pages{101}{120}
\foldertitle{Complexes cubiques : notes manuscrites (s.d.).}
\dating{[à partir de 1976]}
\watermark{Édition de démonstration}

\begin{document}

\page{101}
\note{la page reprend une liste de données dont le début (et l'axiome « (Ka) plus haut ») se trouve avant ce lot}

a) $t_0$ ens.\ de card $3$

b) $\widetilde{\Gamma}$ ens.\ \add{$\neq\emptyset$} avec involution
$\sigma_{\widetilde{\Gamma}}$ (sans pts fixes)

c) Une famille $(\widetilde{E}_\alpha)_{\alpha\in t_0}$ de quotients de
$\widetilde{\Gamma}$ \struck{plus}

Avec les conditions suivantes

$\alpha$) Les $\widetilde{E}_\alpha$ stables par $\sigma_{\widetilde{\Gamma}}$
(dont $\sigma_\alpha$ induites) et les $\sigma_\alpha$ sont sans pts fixes

$\beta_{\widetilde{\Gamma}}$) \struck{$\widetilde{E}$} $\forall\alpha,\beta\in t_0$,
$\alpha\neq\beta$, l'application
\[
  \widetilde{\Gamma}\xrightarrow{\ \widetilde{\varphi}_{\alpha\beta}\ }
  \widetilde{E}_\alpha\times\widetilde{E}_\beta
\]
est \emph{injection}, et si $\widetilde{\Gamma}_{\alpha\beta}$ est son image,
\struck{on a la condition et si $s,s'\in\widetilde{E}_\alpha$ distincts et
conjugués par} \add{est une orbite sous} $\sigma_\alpha$ \struck{($s\neq s'$)} et
$t\in\widetilde{E}_\beta$, on a : $\exists!\ s\in u$ tel que
$(s,t)\in\widetilde{\Gamma}_{\alpha\beta}$
\note{la ligne se lit « si $u\subset\widetilde{E}_\alpha$ est une orbite sous $\sigma_\alpha$ » une fois l'insertion mise en place ; le « ! » est entouré d'un cadre}

($\gamma$) \note{entouré} l'axiome (Ka) plus haut sur $(s_\alpha)_{\alpha\in t_0}$
et $(t_\alpha)_{\alpha\in t_0}$

\struck{On peut déduire de $\beta$), prenant
$\Gamma=\widetilde{\Gamma}/\sigma_{\widetilde{\Gamma}}$,}
\struck{$E_\alpha=\widetilde{E}_\alpha/\sigma_\alpha$, que}
\note{suit un petit diagramme biffé, flèche marquée $\varphi_{\alpha\beta}$}

Soit $\Gamma=\widetilde{\Gamma}/\sigma_{\widetilde{\Gamma}}$,
$E_\alpha=\widetilde{E}_\alpha/\sigma_\alpha$, d'où \uncertain{i.e.}\ des quotients
$E_\alpha$ de $\Gamma$
\[
  \Gamma\xrightarrow{\ \varphi_\alpha\ }E_\alpha ,\qquad
  \Gamma\xrightarrow{\ \varphi\ }\prod_{\alpha\in t_0}E_\alpha
\]
et l'axiome \uncertain{c)} implique que \add{$\forall\alpha$,} $\widetilde{\Gamma}$
s'identifie à l'image inverse des revêtements d'ordre $2$ $\widetilde{E}_\alpha$ de
$E_\alpha$. Donc la famille de données

\page{102}
a) b) c) avec l'axiome $\alpha$ équivaut à : la donnée de l'ens $\Gamma$
\add{$\neq\emptyset$} et de la famille $(E_\alpha)_{\alpha\in t_0}$ de ses quotients
\[
  \varphi_\alpha:\Gamma\to E_\alpha ,\qquad
  \varphi:\Gamma\to\prod_{\alpha\in t_0}E_\alpha ,
\]
\struck{et} \add{puis} des rev.\ d'ordre $2$
\[
  \widetilde{E}_\alpha\xrightarrow{\ p_\alpha\ }E_\alpha\qquad(\alpha\in t_0)
\]
et d'un \emph{syst.\ transitif} d'isom.\ \add{$\Gamma$-} entre les revêtements
$\varphi_\alpha^*(\widetilde{E}_\alpha/E_\alpha)$ de $\Gamma$ (permettant de définir
$\widetilde{\Gamma}$ par identification …) Ceci posé, l'axiome
$\beta_{\widetilde{\Gamma}}$ \struck{implique} \add{équivaut à} l'axiome

$\beta_\Gamma$) Si $\{\alpha,\beta\}\in\mathfrak{P}_2(t_0)$, alors
\[
  \Gamma\xrightarrow{\ \varphi_{\alpha\beta}\ }E_\alpha\times E_\beta
\]
est \emph{bijectif}

La classification des cas « complexes cubiques généralisés » \struck{\ill{}} se
décompose en deux questions exactement

I) Description des systèmes
\[
  \bigl(\Gamma\xrightarrow{\ \varphi_\alpha\ }\Gamma_\alpha\bigr)_{\alpha\in t_0}
\]
de quotients (\uncertain{paramétrés} par $t_0$ de card $3$) d'un ens
$\Gamma\neq\emptyset$, satisfaisant l'axiome $\beta_\Gamma$

\page{103}
II) Pour un tel syst.\ donné, classifier les « rev.\ d'ordre $2$ » des diagrammes
\[
  \bigl(\Gamma\xrightarrow{\ \varphi_\alpha\ }\Gamma_\alpha\bigr)_{\alpha\in t_0}
\]
i.e.\ \struck{des} de ces diagrammes \add{particuliers} de la forme

\begin{tikzcd}[column sep=small, row sep=small]
E_a & & E_c \\
 & \Gamma \arrow[ul, "\varphi_a"] \arrow[ur, "\varphi_c"'] \arrow[d, "\varphi_b"] & \\
 & E_b &
\end{tikzcd}

\note{les indices du diagramme passent de $\alpha$ à $a,b,c$ ; lecture de $E_a$ et $E_b$ incertaine}

puis ceux qui satisfont l'axiome $(\gamma)=(\mathrm{Ka})$
\note{trois traits soulignent la fin de la ligne ; le reste de la page est blanc}

\page{104}
\note{page de calculs, barrée de plusieurs traits obliques}

$s\in\widetilde{E}_a$ définit
$\widetilde{E}_b\overset{f_{s,b}}{\simeq}E_b\times\{\pm1\}$,
$\widetilde{E}_c\overset{f_{s,c}}{\simeq}E_c\times\{\pm1\}$
\qquad $\widetilde{E}_b\simeq E$

\drawing{un triangle de sommets $s,t,u$, traversé par une droite ; à côté \struck{$(s,s',s'_1)$ $4p^2$ $2p^2$} $(s,s'_1)$ ; puis deux segments $s\,t$ et $s'\,t'$}

$2+2c$ \qquad \struck{$2+2c$}

$\mathrm{Card}\,S=\struck{\ill{}}\ 3+6c=3(1+2c)$

$\mathrm{Card}\,E_S(s)=2+2c=2(1+c)$

$\mathrm{Card}\,A=3(1+c)(1+2c)=3(1+3c+2c^2)=3+6c+3c+3(2c^2)$

$\mathrm{Card}\,T=\tfrac13\,\mathrm{Card}\,A=(1+c)(1+2c)=1+3c+2c^2$

\note{diagramme biffé :
$\widetilde{\Gamma}\to\widetilde{E}_\alpha\times\widetilde{E}_\beta\to\widetilde{E}_\alpha\times E_\beta\to\widetilde{E}$,
avec $\widetilde{\Gamma}\to\Gamma\simeq E_\alpha\times E_\beta$ et une flèche vers $E$}

$\mathrm{Card}\,T(s)=1+c$ \struck{\ill{}} \qquad
\struck{$\mathrm{Card}\,T(s)=1+c$}

$\mathrm{Card}\,T=\tfrac13\,\mathrm{Card}\,S\,(1+c)$

$\mathrm{Card}\,\widetilde{\Gamma}=2c^2$, \quad $\mathrm{Card}\,\Gamma=c^2$,
\quad $\mathrm{Card}\,\widetilde{E}_a=2c$, \quad $\mathrm{Card}\,E_a=c$

\drawing{petit graphe : un rectangle partagé par une barre médiane, les deux côtés verticaux doublés d'arcs ; quelques points isolés}

\page{105}
\drawing{en tête de page : une étoile de triangles (deux sommets épais reliés, portant chacun plusieurs triangles), marquée $\sigma_{t_0}$ ; un signe en gras ; deux triangles opposés par un sommet ; un triangle isolé}

\ill{} \note{mot entouré} \ill{} : $0$ \uncertain{boucles} a) — des points isolés
— \uncertain{pas un sommet porte $0$ triangle} (\ill{} : $0$ \uncertain{arête})

\ill{} \note{mot entouré} $n=1$ \uncertain{boucle} b) — un triangle
— pas un sommet porte $1$ triangle exactement (ou $2$ arêtes
\uncertain{contiguës})

c) $n\geqslant2$ \uncertain{boucles} — une étoile de quatre triangles en un sommet
— \struck{\ill{}} Chaînes de triangles \ill{} \struck{à deux sommets} tous les
sommets \ill{} sauf $1$ d'ordre $2$, \struck{l'autre} (l'exceptionnel est d'ordre
$\geqslant4$)

pour tous les sommets sauf $1$ porte exactement un triangle, pour l'excep.\ il y en
a au moins deux [Alors il y a au moins $4$ de sommets par lesquels passe
exactement $1$ triangle …]

\marginal{Car il existe un autre (\uncertain{par} \ill{}) auquel \ill{} autre \ill{}
(\uncertain{lié}) — unique \ill{} \ill{} triangles)}

\ill{}, Si un sommet porte au moins $2$ triangles
\note{la page s'arrête là ; la suite manque}

\section{Automorphismes des graphes trialitaires}
\note{titre de sa main, souligné ; « complexes » y est biffé et remplacé par « graphes »}

\page{106}
Soit $(S,A)$ un \struck{complexe} \add{graphe} trialitaire,
$\underline{t}=\{a,b,c\}$ un triangle. On a défini une bijection involution
\struck{$\sigma_{\underline{t}}$ de $S$, …}
\[
  \sigma_{\underline{t}}=\sigma_{\{a,b,c\}}:S\to S
\]
par
\[
  \sigma_{\underline{t}}\,|\,\underline{t}=\mathrm{id}_{\underline{t}},\qquad
  \sigma_{\underline{t}}\,|\,E_a^{\underline{t}}=\sigma_a^{\underline{t}}
\]
(transforme $x\in E_a$ en l'unique él.\ de $E_a$ lié à $x$)

\emph{Proposition} a) $\sigma_{\underline{t}}$ est un automorphisme du
\struck{complexe} graphe $(S,A)$, et $\underline{t}$ est l'ens des pts fixes de
$\sigma_{\underline{t}}$.

b) L'application $\underline{t}\mapsto\sigma_{\underline{t}}$ de $T$ dans
$\mathcal{G}=\mathrm{Aut}(S,A)$ est injective, et a comme image l'ens des autom.\
involutifs \add{$u$} de $(S,A)$ tels que \struck{l'ens des}
$\mathrm{card}\,S^u=3$

C'est évident…

Soit $\mathcal{G}=\mathrm{Aut}(S,A)$. Les $\sigma_{\underline{t}}$ sont
$\in\mathcal{G}$. Soit $\mathcal{G}^{\circ}$ le groupe qu'ils engendrent. Si
$g\in\mathcal{G}$, on a
\[
  \mathrm{int}(g).\sigma_{\underline{t}}=g\,\sigma_{\underline{t}}\,g^{-1}
  =\sigma_{g(\underline{t})}
\]
donc $\mathcal{G}^{\circ}$ est un ss-groupe invariant.

\page{107}
\emph{Théorème} $\mathcal{G}^{\circ}$ est transitif sur l'ens des arêtes de
$(S,A)$ (i.e.\ des couples $(s,t)$ de deux pts liés) et
\add{\uncertain{l'ens des faces} \ill{}} a fortiori sur l'ens des sommets,
\struck{et} sur l'ens des triangles \add{ensemble} (qui \ill{} \ill{} : des ens.\
\uncertain{partiels} de l'ens \ill{} …) \note{la fin de l'énoncé est perdue dans les insertions ; on lit encore « dans $\mathcal{G}^{\circ}$ (a fortiori dans $\mathcal{G}$) » et « \ill{} ici deux.) »}

\marginal{Sauf exception $c\,(=\mathrm{Card}\,E_a)=1$ i.e.\ $(S,A)$ est le
bi-triangle \ill{}. Alors $\mathcal{G}$ a les transitivités \uncertain{énoncées},
mais $\mathcal{G}^{\circ}\simeq\mathfrak{S}_3\times\mathfrak{S}_3$ \ill{} est
d'indice $2$ dans $\mathcal{G}$, \uncertain{ici}
$\mathcal{G}\simeq\mathfrak{S}_2.(\mathfrak{S}_3\times\mathfrak{S}_3)$ ;
$\mathcal{G}^{\circ}$ est encore transitif sur $S$ \ill{}, mais pas sur $A$ (\ill{})
\ill{} triangles \ill{} les \uncertain{pointillés} …}

\emph{Cor.} Les $\sigma_{\underline{t}}$ ($\underline{t}\in T$) sont conjugués, et
forment exactement une classe de conjugaison d'éléments de $\mathcal{G}^{\circ}$
(ou de $\mathcal{G}$)

Résulte des deux lemmes suivants

\emph{Lemme 1} Soient $s,t\in S$, alors

a) Si $s,t$ sont liés, $\exists\,\underline{t}\in T$ tel que
$t=\sigma_{\underline{t}}\,s$

b) \struck{Dans tous les cas, $s,t$ étant non liés} $\exists\,u\in S$ tel que $s$
lié à $u$ et $u$ lié à $t$, donc (lemme a) des $\underline{t},\underline{t}'\in T$
tels que $u=\sigma_{\underline{t}}\,s$, $t=\sigma_{\underline{t}'}\,u$, donc
$t=\sigma_{\underline{t}'}\sigma_{\underline{t}}\,s$.

Le lemme a) est clair (\struck{prendre pour} \add{\uncertain{savoir}} \ill{}
sommet des triangles définis par $\{s,t\}$, \struck{le triangle défini \ill{} un
triangle}, et \ill{} distinct des \uncertain{trois} triangles
\uncertain{précédents}) ; b) \ill{} \ill{} au moins $3$ triangles incidents en
chaque sommet … (\uncertain{Supposons} $c\neq1$ i.e.\ \ill{})

\emph{Lemme 2} Soient $a\in S$, \struck{\ill{}} et $s,t$ deux él.\ de $S$ liés à
$a$ \add{\ill{}}. Alors $\exists\,g\in\mathcal{G}^{\circ}$ (produit de deux
$\sigma_{\underline{t}}\sigma_{\underline{t}'}$, où $\underline{t},\underline{t}'$
$\mapsto$ $2$ triangles \uncertain{incidents convenables}) tels que $u(a)=a$,
$u(s)=t$, $u(t)=s$ \note{la page passe de $g$ à $u$ pour le même élément}

\page{108}
Plus précisément, si $\underline{t}_0$ \add{$=\{a,b,c\}$} est un triangle de
sommet $a$ ne contenant ni $s$ ni $t$ (il y en a par l'hyp.\ que $c\neq1$ i.e.\
$c\geqslant2$) on peut choisir de plus $u$ \uncertain{lié à}
$\underline{t},\underline{t}'$) tel que
\struck{$u\,|\,\underline{t}_0=\mathrm{id}$ \ill{}}
\add{$u\,|\,\underline{t}_0=\mathrm{id}$ i.e.}
\[
  u\,a=a,\quad u\,b=b,\quad u\,c=c
\]
\struck{\ill{}, $u\,a=a$, $u\,b=c$, $u\,c=b$.}

\struck{On peut trouver $w\in\mathcal{G}_0$ tel que
$w\,|\,\underline{t}_0=\underline{t}_0$}
\struck{\ill{} $w\,a=a$, $w\,b=b$, $w\,c=c$) et $w\,s=t$}

\drawing{en marge gauche, deux figures de triangles accolés : la première avec les triangles $\underline{t}$, $\underline{t}_1$, $\underline{t}_0$ et les sommets $s,s',t,t',a,b,c,u,u',v,v'$ ; la seconde avec $t,t',b,c,a,s,u,\bar u$}

\note{la suite de la page est barrée de longs traits obliques}

\emph{Recette} On \struck{choisit $u$ lié à $t$ et $s$} \add{complète les
triangles} $\underline{t}=\{a,s,s'\}$ et \struck{\ill{}}
$\underline{t}_0=\{a,b,c\}$ (\uncertain{pas} $s'$ et $c$) (ce qui \ill{} …) on
choisit $\widetilde{E}_a,\widetilde{E}_b,\widetilde{E}_c$ etc …) on choisit $u$
\note{entouré} \add{$\neq a$} lié à $b$ et à $s$ i.e.\ $u\in\widetilde{E}_b$ lié
à $s$, on complète en triangle $\{s,u,v\}$ un $(a,s,c)$ donc
\[
  u\in\widetilde{E}_b,\ v\in\widetilde{E}_c \qquad
  \{s,u,v\}\in\widetilde{T}\subset T
\]
\[
  u'\in\widetilde{E}_b,\ v'\in\widetilde{E}_c \qquad
  \{s',u',v'\}\in\widetilde{T}\subset T
\]
\[
  [\,(b,u,u')\in T,\quad(c,v,v')\in T\,]
\]
$t,t'$, \ill{} $S$

\struck{On peut} choisit \ill{} : \uncertain{échange} ; $t$ lié à $u$,
\struck{\ill{}} complète dans un triangle
\[
  (t,u,v_1)\in\widetilde{\Gamma}
\]
Alors $t'$ est lié ($\widetilde{E}_c$) à $v'$, donc par à $v'$, donc à $v$ … de
complète
\[
  (t',v,u_1)\in\widetilde{\Gamma}\subset T\quad(\widetilde{E}_b)
\]
Et comme $u_1,v_1$ — a aussi $u'_1\in\widetilde{E}_b$, $v'_1\in\widetilde{E}_c$,
$\{b,u_1,u'_1\},\{c,v_1,v'_1\}\in T$

\page{109}
\struck{\ill{}, on prendra
$\underline{\tau}=\{t',u'_1,v'_1\}$, $\underline{\tau}'=\{t,u_1,v_1\}$}
\note{ces deux lignes sont barrées de traits obliques}

Changeons de notation donnés $a,s,t$ ($s,t$ liés à $a$, distincts et non liés) on
complète triangles (distincts) $\{a,s,s'\}$, $\{a,t,t'\}$, on a triangle
$\underline{t}_0$ de sommet $a$ distinct des $2$ précédents. On choisit $u\neq a$
lié à $s$ et à $t$ (il est lié à un unique sommet de $\underline{t}_0$, qui est
distinct de $a$, disons $b$), on prend $\{s,u,v\}\in T$, $\{t,u,w\}\in T$,
\struck{il} \add{\uncertain{qui}} sont liés à $c$ (et sont distincts de $a,s$), on
prend les triangles $\underline{\sigma}=\{c,v,v'\}$ et
$\underline{\tau}=\{c,w,w'\}$, et on considère les symétries
$\sigma_{\underline{\sigma}}$, $\sigma_{\underline{\tau}}$, qui \uncertain{commutent}
… (\ill{} triangles \uncertain{sécants}). Car $\underline{\sigma},\underline{\tau}$
sont des \ill{}

On a
\[
  \sigma_{\underline{\sigma}}\,s=u \quad\text{car } \{s,u,v\}\in T,\ v\in\underline{\sigma}
\]
\[
  \sigma_{\underline{\tau}}\,u=t \quad\text{car } \{t,u,w\}\in T,\ w\in\underline{\tau}
\]
et fixent $c$ \note{la lettre $\underline{\sigma}$ nommant un triangle est d'une lecture incertaine}

enfin, $\sigma_{\underline{\sigma}}$ et $\sigma_{\underline{\tau}}$ échangent $a$
et $b$, donc leur produit fixe les pts $a,b,c$. Donc
\[
  W=\sigma_{\underline{\tau}}\,\sigma_{\underline{\sigma}}
\]
transforme $(a,s)$ en $(a,t)$ et induit l'identité dans $(a,b,c)$

\drawing{en marge gauche, deux figures : les triangles $\{a,s,s'\}$, $\{a,t,t'\}$ autour de $a$ ; puis le triangle $a,b,c$ avec $u$ et les triangles $\{c,v,v'\}$, $\{c,w,w'\}$}

\marginal{\ill{} \ill{} \ill{} … $\sigma_{\underline{\tau}}\,u=s$,
$\sigma_{\underline{\sigma}}\,s=u$ donc \ill{} $u=\mathrm{id}$, et transforme
$\{s,u,v\}\in\widetilde{\Gamma}$, $\{t,u,w\}\in\widetilde{\Gamma}$ \ill{}}

\struck{\emph{NB} Supposons qu'il existe $r\in\widetilde{E}_c$ lié à $s$ et $t$.
Comme $s$ lié à $v$ (donc \ill{} à $v'$), $t$ lié à $w$ (donc \ill{} à $w'$) …
$r\neq v,w$}
\note{ce NB est barré de traits obliques ; il se poursuit en tête de la p.~112, elle aussi barrée}

\page{110}
\emph{Corollaire} Considérons \struck{le} le groupe \add{(invariant)
$\mathcal{H}^{\circ}$ des} \struck{\ill{}}
$\mathcal{H}=\mathrm{Aut}(S,A;a,b,c)$ \struck{engendré} \add{fourni} par les
produits d'un nb pair de symétries $\sigma_{\underline{t}}$ relatifs à des
triangles verticaux. Alors $\mathcal{H}^{\circ}$ est transitif sur l'ens
$\widetilde{\Gamma}$ des triangles supérieurs. De façon précise, si
\struck{\ill{}} $\underline{s},\underline{t}\in\widetilde{\Gamma}$, et si
$a\in\underline{t}_0$, soit \struck{\ill{}} $b\in t_0$, $b\neq a$ tel que
$\exists\,u\in\widetilde{E}_b$ lié aux deux sommets \add{$s,t$} de
$\underline{s}$, $\underline{t}$ sur $a$, alors $\exists$ \struck{produit} $U$ de
deux symétries \struck{relatifs} $\sigma_v,\sigma_w$
($v,w\in E_c=\widetilde{E}_c/2$) tel que $U\,s=t$, \uncertain{donc}
$U\underline{s}$ et $\underline{t}$ coïncident \emph{en $a$} ; et \ill{} \ill{}
sommets \ill{} \emph{\uncertain{choix}} \ill{} trouver $V$, produit de deux
symétries \uncertain{relatives} en $b$, en $c$, tel que
\[
  V\,U\,\underline{s}=\underline{t}.
\]
\struck{Si on \ill{} pour \ill{},}
\struck{$V$ et $U$ commutent, $VU$ est un automorphisme involutif}
\note{la moitié inférieure de ce corollaire est barrée d'un grand trait ondulé ; le reste de la page est blanc}

\page{111}
\emph{Lemme 1} (\uncertain{lors} $c$ \uncertain{fini} $\neq1,2$) Soient
$s,t\in\widetilde{E}_a$, alors $\exists\,s_1\in\widetilde{E}_a$ et
$u,v\in\widetilde{E}_b$ tels que $u$ soit lié à $s$ et $s_1$, $v$ lié à $s_1$ et
$t$.

\marginal{\uncertain{c'est faux} si $c=1,2$}

Soit $\widetilde{E}_b(s)=\{u\in\widetilde{E}_b\mid u\text{ lié à }s\}$, de même
$\widetilde{E}_b(t)$, ils sont de cardinal $c$ (des sections de
$\widetilde{E}_b$ sur $E_b$). Si
$\widetilde{E}_b(s)\cap\widetilde{E}_b(t)\neq\emptyset$, il suffit de prendre
$u=v$ \uncertain{élt} dans cette intersection et $s_1=s$. Nous pouvons donc
supposer
\[
  \widetilde{E}_b=\underset{B'}{\widetilde{E}_b(s)}\amalg
  \underset{B''}{\widetilde{E}_b(t)}
\]
(comme \uncertain{même} $\widetilde{E}_b(t)=\widetilde{E}_b(s')$ \ill{} cela
impliquera $t=s'$, mais \ill{} \ill{} \ill{}) On cherche \struck{\ill{}} $\bar s_1$
tel que
\[
  \widetilde{E}_b(\bar s_1)\cap\widetilde{E}_b(s)\neq\emptyset,\qquad
  E_b(\bar s_1)\cap\widetilde{E}_b(t)\neq\emptyset .
\]
Si un tel $\bar s$ n'existait pas, on aurait
$\forall\,\bar r\in\widetilde{E}_a$, $\widetilde{E}_b(\bar r)$ est contenu dans
\struck{$\widetilde{E}_b(s)$} $B'$ ou $B''$ et pour des raisons de cardinal
(\uncertain{comme il en est} \ill{} de cardinal $c$) on aurait
\[
  \forall\,r\in A=\widetilde{E}_a,\ \text{\uncertain{ou}}\quad
  \widetilde{E}_b(r)=B'\ \text{ou}\ \widetilde{E}_b(r)=B''
\]
Soit
\[
  A'=\{r\in A\mid\widetilde{E}_b(r)=B'\},\qquad
  A''=\{r\in A\mid\widetilde{E}_b(r)=B''\}=\complement_A A'
\]
donc \struck{$A'=\sigma_{\underline{t}_0}A''$, $A''=\sigma_{b}A'$, et \ill{}}
\ill{} $r,r'$ de \ill{} de $A$ liés, un et un seul de $r,r'$ est dans $A'$,
l'autre dans $A''$ — donc $A',A''$ \ill{} en deux \uncertain{parties}
complémentaires de $\widetilde{E}_a$ \ill{}

Soit \struck{$r_1,r_2\in\widetilde{E}_a$, \ill{}} $r_1\neq r_2$
\struck{$r_1\neq r'_1$} $r_1\in A'$, $r'_2\in A''$
\struck{\ill{}}

Soit \struck{\ill{}} $r_1,r_2\in A'$, \struck{$r_2\in A''$ \ill{}} et
$r'_2\neq(r_1)'$, \ill{} $\neq$, dans $A''$, \ill{} $\widetilde{E}_b$ lié :
$r_1,r'_2$, $2$ fois, dans l'ens des pts de $\widetilde{E}_c$ liés : $r_1$ et
$r'_2$ \ill{} $2$ fois \ill{}

\page{112}
\note{page barrée de deux longs traits obliques ; elle continue le NB biffé de la p.~109}

et communs sont liés : \ill{} \ill{} $\{s,u,v\}$ et $\{s,u,w\}$ \struck{\ill{}}
des triangles), \uncertain{mais} $t$ \ill{} $u$, \uncertain{donc} \ill{}
$r\neq u,w$. Donc $\sigma_{\underline{s}}$ et $\sigma_{\underline{t}}$
\uncertain{sont} l'identité sur $r$, donc \struck{\ill{}} \ill{}
$\sigma_{\underline{s}}\sigma_{\underline{t}}$ \ill{}

Donc, si \struck{\ill{}} $\underline{s},\underline{t}\in\widetilde{\Gamma}$ sont
deux triangles \uncertain{supérieurs} qui coïncident \struck{en $a$} au sommet $c$
(\uncertain{sur} $c$) et qui ont \struck{\ill{}} les sommets $s,t$ au-dessus de
$a$, \uncertain{alors}, si on peut trouver $u\in\widetilde{E}_b$ lié \ill{} fois :
$s$ et $t$, alors

\drawing{un triangle de sommets $a,b,c$ ; au-dessus, les points $s$ et $t$ reliés à un point marqué ; $u$ à droite}

\page{113}
de cardinal $c$, donc par des raisons de cardinal \ill{} aussi
\[
  \widetilde{E}_c(r_1)=\widetilde{E}_c(r'_2)
\]
\note{ce début enchaîne sur la fin de la p.~111}

Si \uncertain{maintenant} $c\geqslant3$ \add{$r_1\neq r_2$}, de sorte que pour
$r_1,r_2\in A'$ distincts, $\exists\,r_3\in A'$ tel que $r_3\neq r_1,r_2$. On trouve
donc
\[
  \widetilde{E}_c(r_1)=\widetilde{E}_c(r_2)
\]
\add{$=\widetilde{E}_c(r'_3)$},
donc tous les $\widetilde{E}_c(r_i)$ \struck{($r_i\in A'$)}
\struck{sont égaux, et $\widetilde{E}_c(r_i)$ sont égaux}
sont égaux et égaux aux $\widetilde{E}_c(r'_i)$, or
$\widetilde{E}_c(r'_i)=\complement\,\widetilde{E}_c(r_i)$, ce qui est absurde.

\emph{Lemme 2} Soient \struck{$v_1,v_2\in\widetilde{E}_b$,
$u\in\widetilde{E}_b$ lié à $s$ et $t$,}
$\underline{t}_1=(u,v_1,w_1)$, \add{$\underline{t}_2=$} $(u,v_2,w_2)$ deux
triangles sup.\ sur $(a,b,c)$, ayant le sommet $u$ (sur $b$) en commun.
Considérons $\pi=\sigma_{w_1}\sigma_{v_2}$ \ill{} \add{\uncertain{l'autre}} \ill{}
qui fixe $a,b,c$) alors $\pi$ \add{\uncertain{échange} $w_1$ et $w_2$}
\struck{échange les deux triangles parallèles, i.e.\ $\pi u=u$,
$\pi v_1=w_2$, $\pi w_1=w_2$} (si $c\neq1,2$)
\note{les indices de $\pi$ sont surchargés ; lecture de $\sigma_{w_1}\sigma_{v_2}$ incertaine}

\emph{Corollaire} \struck{Soient} $a,b$ deux sommets distincts de
$\underline{t}_0$, $\mathcal{G}_b$ le s/groupe \add{des \uncertain{autom}} des
automorphismes de $(S,A)$ engendré par les produits de \struck{deux}
\struck{symétries} par $\sigma_{v_1}\sigma_{v_2}$, avec
$v_1,v_2\in\widetilde{E}_b$ \ill{} (donc $\mathcal{G}_b$ est un
\struck{\ill{}} $\mathbb{F}_2$-vectoriel \ill{}, et il respecte $a,b,c$) Alors
$\mathcal{G}_b$ est transitif sur $\widetilde{E}_a$, de façon plus précise, si

\page{114}
$s,t\in\widetilde{E}_a$, $\exists\,\pi\in\mathcal{G}_b$ produit de $4$
\struck{\ill{}} symétries de $2$ \uncertain{formes} $\sigma_v$
($v\in E_b$) tel que $t=\pi s$.

\emph{Corollaire plus fort} $\mathcal{G}_b$ est transitif sur
$\widetilde{\Gamma}$, plus précisément, si $\bar s,\bar t\in\widetilde{\Gamma}$, il
existe un $\pi\in\mathcal{G}_b$ produit de $6$ symétries de \uncertain{la forme}
$\sigma_w$ ($w\in E_b$) tel que $\bar t=\pi\bar s$.
\note{ce « corollaire plus fort » est barré de quatre traits obliques ; le reste de la page est blanc}

\page{115}
nb de triangles épinglés $=$ nb d'arêtes
\[
  \mathrm{trep}=6(1+2c)(1+c)=6(1+3c+2c^2)
\]
nb de bitriangles épinglés : \struck{\ill{}}
\[
  \mathrm{bitrep}=2c^2\,\mathrm{trep}=12c^2(1+2c)(1+c)
\]
\[
  =12c^2(1+3c+2c^2)
\]

\emph{NB} Si $c=2$, \struck{les} automorphismes de $S=(S,A)$ qui
\struck{\ill{}} \add{\uncertain{induisent}} sur un bitriangle \emph{épinglé} est
l'identité. Donc le groupe $\mathcal{G}$ des automorphismes \struck{\ill{}} est
$=$ le groupe \struck{\ill{}} des automorphismes « intérieurs » i.e.\ le s-gr de
$\mathcal{G}$ engendré par les $\sigma_{\underline{t}}$, $\underline{t}\in T$) est
\emph{simplement transitif} sur l'ens des bitriangles épinglés. Donc ici
\[
  \mathcal{G}=\mathcal{G}^{\circ}\quad\text{d'ordre}\quad 12c^2(1+2c)(1+c)
\]
\[
  =12.4.5.3=2^4.3^2.5=720
\]
\note{la formule est surchargée de corrections des facteurs ; le calcul $12.4.5.3$ porte aussi un chiffre biffé}

\emph{NB} Si $c=1$,
$\mathcal{G}^{\circ}\simeq\mathfrak{S}_3\times\mathfrak{S}_3$ \struck{\ill{}} est
d'indice $2$ dans
\[
  \mathcal{G}\simeq\mathfrak{S}_2.(\mathfrak{S}_3\times\mathfrak{S}_3),
\]
qui est d'ordre $72=2^3 3^2$. On l'obtient aussi comme le nb des bitriangles
épinglés (puisque $\mathcal{G}$ est simpl.\ transitif dessus) qui est
\[
  12c^2(1+2c)(1+c)=12.1.3.2=72
\]
$\Downarrow$

Si $c\geqslant1$, le nb de bitriangles $=\frac1{72}$ nb des bitriangles épinglés
\[
  =\frac16\,c^2(1+2c)(1+c)
\]
\note{en bas à gauche, quelques chiffres de brouillon, en partie biffés ($720$…)}

\page{116}
\note{page barrée d'un long trait oblique, qui ne semble pas annuler les formules}

\[
  xyz=1
\]
\drawing{un triangle aux sommets numérotés $1,2,3$}
\[
  \varphi_a:(x,y,z)\longmapsto(x,\,a^{-1}z^{-1},\,a^{-1}y^{-1})
\]
\marginal{$a^{-1}=yz$, \quad $y=a^{-1}z^{-1}$, \quad $z=y^{-1}a^{-1}$}
\[
  xyz=1\Longrightarrow x\,a^{-1}z^{-1}a^{-1}y^{-1}=1
\]
\uncertain{jamais} $xyz=1$, — \uncertain{donne} $a^{-2}=1$ \uncertain{donc}
$a^2=1$ : \uncertain{tout élément} est d'ordre $2$, ce qui implique que $G$ est
\struck{$abab=1$} commutatif, donc un vectoriel sur $\mathbb{F}_2$ — et dans la
condition \uncertain{est} \uncertain{satisfaite}.
\[
  \varphi_c\quad(x,y,z)\longmapsto(x,\,z+c,\,y+c)
\]
\[
  \varphi_a\varphi_b:\ (x,y,z)\longmapsto(x,\,y+c,\,z+c)\qquad c=a+b
\]

$\mathcal{G}$ opère transitivement sur $\widetilde{\Gamma}$

\begin{tikzcd}[column sep=small, row sep=small]
 & \mathcal{G}_a \arrow[d, hook] & \\
\mathcal{G}_b \arrow[r, hook] & \mathcal{G} & \mathcal{G}_c \arrow[l, hook']
\end{tikzcd}

ss-groupes invariants qui sont des $\mathbb{F}_2$-vectoriels.

$\mathcal{G}$ opère transitivement sur
$\widetilde{\Gamma}\simeq\mathcal{G}/\mathfrak{g}$ ($\mathfrak{g}$ stabilisateur de
$\underline{t}\in\widetilde{\Gamma}$)

\begin{tikzcd}[column sep=small, row sep=small]
 & \widetilde{E}_a\simeq\widetilde{E} & \\
 & \widetilde{\Gamma} \arrow[u] \arrow[dl] \arrow[dr] & \\
\widetilde{E}\simeq\widetilde{E}_b & & \widetilde{E}_c\simeq\widetilde{E}
\end{tikzcd}

$\mathcal{G}_a$, \struck{\ill{}} $\mathcal{G}_b$ opère trans.\ sur
$\widetilde{E}_b$, \add{$\widetilde{E}_c$}), et trivialement sur
$\widetilde{E}_a$

$\mathcal{G}_a$ isom.\ au groupe des \uncertain{parties} paires de $E_a$
\[
  [\,\mathcal{G}_a\to\mathrm{Aut}_{E_a}(\widetilde{E}_a)\ \text{est injectif}\,]
\]

\page{117}
\[
  \alpha,\beta,\gamma\in[-1,2]\qquad\alpha\leqslant\beta\leqslant\gamma
\]
$c_{\beta\gamma}=$ Nb de structures $\tau_\beta$ dans un $\tau_\gamma$

$c'_{\beta\gamma}=$ —— épinglées $\tau_\beta$ dans un $\tau_\gamma$

$N_\beta=$ \add{$c'_{\beta\beta}=$} ordre groupe autom de $\tau_\beta$
\[
  (1)\quad \boxed{c'_{\beta\gamma}=N_\beta\,c_{\beta\gamma}}\qquad\text{i.e.}\qquad
  \boxed{c_{\beta\gamma}=\frac{c'_{\beta\gamma}}{c'_{\beta\beta}}}
\]
\note{au-dessus du second cadre, biffé : « $N_\beta$ … $N_\alpha=c'_{\alpha\alpha}$ »}
\[
  (2)\quad \boxed{c'_{\beta\gamma}=c'_{\alpha,\gamma}\,\nu_{\alpha,\beta,\gamma}}
\]
$\nu_{\alpha\beta\gamma}=$ nb des structures épinglées $\tau_\beta$ dans un
$\tau_\gamma$, prolongeant une structure épinglée $\tau_\alpha$ \add{donnée},

On calcule les $c_{\beta\gamma}$, $c'_{\beta\gamma}$ ($\gamma$ fixé,
\add{tous les} $\beta\leqslant\gamma$) par récurrence sur $\beta$,
\struck{\ill{}} \uncertain{à partir} de $c'_{\beta\gamma}$ \ill{} \ill{} de
\ill{} $N_\beta$ \ill{}
\[
  c'_{\beta\gamma}=N_\beta\,c_{\beta\gamma}
\]
Si $\beta=\gamma$, $c_{\beta\gamma}=1$, \struck{$c'_{\beta\gamma}$} il suffit de
connaître $c'_{\beta\gamma}$. \struck{Et choisissant} et pour $\beta=-1$
($\tau_\beta$ triangle) on a la formule
\[
  c_{-1,\gamma}=\text{nb de triangles \add{de $\tau_\gamma$}}
  =(1+c_\gamma)(1+2c_\gamma)
\]
\[
  c'_{-1,\gamma}=\text{nb de triangles épinglés de }\tau_\gamma
  =6(1+c_\gamma)(1+2c_\gamma)
\]
On \uncertain{récurre} sur $\beta$ \add{(en prenant $\alpha=\beta-1$)}
\struck{On a donc} si $\beta\neq-1$ (donc $\gamma\geqslant0$, i.e.\
$\gamma=0,1$ ou $2$) \add{par $\alpha<\beta$} alors la formule (2)
\uncertain{nous} donne $c'_{\beta\gamma}$ en fonction de $c'_{\alpha\gamma}$
(déjà connu) et $\nu_{\alpha\beta\gamma}$ (qu'il faut encore déterminer).
\struck{Reste (\ill{} \ill{} $\nu_{\alpha\beta\gamma}$) : \ill{} \ill{}
$N_\gamma$.}

\note{la fin de la page est barrée de traits obliques}

On \struck{\ill{}} peut prendre $\beta=\gamma-1$, \uncertain{considérer} le
groupe $\mathcal{G}_\gamma$ opérant sur l'ens des $\tau_\beta$
\uncertain{épinglés} (au nombre de $c'_{\beta\gamma}$) dans le $\tau_\gamma$,
\ill{} $\mathfrak{g}_{\beta\gamma}$ \uncertain{stabilisateur} \ill{} il est
d'indice $c'_{\beta\gamma}$, et son ordre est $N_{\beta\gamma}$ \ill{} groupe des

\page{118}
automorphismes de $\tau_\gamma$ induisant l'identité sur $\simeq\tau_\beta$
plongé, d'où
\[
  N_\gamma=N_\beta\,N_{\beta\gamma}
\]
Donc il suffit de déterminer les \add{si \struck{$0$} $1\leqslant\beta<\gamma$}
\[
  \nu_{\beta-1,\beta,\gamma}\quad\text{et}\quad
  N_{\gamma-1,\gamma}\ \text{pour}\ \gamma\geqslant0
\]
\note{devant $\nu_{\beta-1,\beta,\gamma}$, un mot biffé illisible}

On trouve
\[
  \nu_{-1,0,\gamma}=2c_\gamma^2\quad
  \begin{cases}\gamma=0 & 2\\ \gamma=1 & 8=2^3\\ \gamma=2 & 32=2^5\end{cases}
\]
\[
  \nu_{0,1,\gamma}=c_\gamma-1\quad
  \begin{cases}\gamma=1 & 1\\ \gamma=2 & 3\end{cases}
\]
\[
  \nu_{1,2,\gamma}=c_\gamma-2=2\ \text{si}\ \gamma=2
\]
\note{plusieurs essais biffés à gauche de ces valeurs, dont $N_{-1,0}$, $N_{0,1}$, $N_{1,2}$ et un « $\nu_{0,1,2}$ » recouvert}

\marginal{bitriangle en bitriangle}

\drawing{à gauche, des esquisses de bitriangles : quatre petits triangles en ligne, puis une figure de triangles accolés}

$c'_{\beta\gamma}$ ($\beta\leqslant\gamma$)

\[
  c'_{-1,\gamma}=6(1+c_\gamma)(1+2c_\gamma)\qquad
  6\cdot\left\{\begin{array}{l}1.1\\2.3\\3.5\\5.9\end{array}\right.
\]

\[
  \begin{array}{c|c|c|c|c}
  \beta\backslash\gamma & -1 & 0 & 1 & 2\\\hline
  -1 & 6=2.3 & 36=2^2.3^2 & 90=2.3^2.5 & 270=2.3^3.5\\\hline
  0 & 0 & 2^3.3^2 & 2^4.3^2.5 & 2^6.3^3.5\\\hline
  1 & 0 & 0 & 2^4.3^2.5 & 2^6.3^4.5\\
    &   &   & =6!=720 & \\\hline
  2 & 0 & 0 & 0 & 2^7.3^4.5
  \end{array}
\]

\page{119}
\[
  c'_{-1,\gamma}=6(1+c_\gamma)(1+2c_\gamma)\qquad\gamma\geqslant-1
\]
\[
  c'_{0,\gamma}=c'_{-1,\gamma}\,\nu_{-1,0,\gamma}
  =12c_\gamma^2(1+c_\gamma)(1+2c_\gamma)\qquad\gamma\geqslant0
\]
\[
  c'_{1,\gamma}=c'_{0,\gamma}\,\nu_{0,1,\gamma}
  =12c_\gamma^2(c_\gamma+1)(c_\gamma-1)(2c_\gamma+1)\qquad\gamma\geqslant1
\]
\[
  c'_{2,\gamma}=c'_{1\gamma}\,\nu_{1,2,\gamma}
  =12c_\gamma^2(c_\gamma+1)(c_\gamma-1)(c_\gamma-2)(2c_\gamma+1)
\]
\[
  \gamma\geqslant2\ \text{i.e.}\ \gamma=2
\]
$c_{-1,\gamma}=$ \add{$\tfrac16\,c'_{-1,\gamma}$} \struck{\ill{}}
$(1+c_\gamma)(1+2c_\gamma)$

\[
  c_{0,\gamma}=\frac1{72}\,c_{0,\gamma}
  =\frac16\,c_\gamma^2(c_\gamma+1)(c_\gamma-1)(2c_\gamma+1)
\]
\note{ainsi sur la page : $c_{0,\gamma}$ sans prime au second membre, et un facteur $(c_\gamma-1)$ que la ligne $c'_{0,\gamma}$ ne porte pas ; le tableau ci-dessous ($10$, $120$) suit la formule sans ce facteur}
\[
  c_{1,\gamma}=\frac1{720}\,c'_{1,\gamma}
  =\frac1{60}\,c_\gamma^2(c_\gamma-1)(c_\gamma+1)(2c_\gamma+1)
\]
\[
  c_{2,\gamma}=1\ \text{—}
\]
\[
  \begin{array}{c|c|c|c|c}
  \beta\backslash\gamma & -1 & 0 & 1 & 2\\\hline
  -1 & 1 & 6=2.3 & 15=3.5 & 45=3^2.5\\\hline
  0 & 0 & 1 & 10=2.5 & 120=2^3.3.5\\\hline
  1 & 0 & 0 & 1 & 36=2^2.3^2\\\hline
  2 & 0 & 0 & 0 & 1
  \end{array}
\]

\page{120}
\note{page écrite en travers de la feuille ; un grand tableau à colonnes $c=0,1,2,4$, transcrit ligne par ligne}

\[
  \begin{array}{l|c|c|c|c}
   & c=0 & 1 & 2 & 4\\\hline
  \text{Nb de triangles incidents à un sommet : } c+1=e & 1 & 2 & 3 & 5\\
  \text{Nb de sommets —— : } 2(c+1) & 2 & 4 & 6 & 10\\
  \text{Nb de sommets : } 3(1+2c)=3+6c & 3 & 9 & 15 & 27
  \end{array}
\]
\[
  \begin{array}{l|c|c|c|c}
  \text{Nb des arêtes : } 3(c+1)(1+2c) & 3 & 18 & 45 & 135\\
  \quad =3(1+3c+2c^2) & & & & \\\hline
  \text{Nb de triangles } (c+1)(1+2c)=1+3c+2c^2 & 1 & 6 & 15 & 45\\
  \text{Nb de triangles verticaux } =3c & 0 & 3 & 6 & 12\\
  \text{\phantom{Nb de triangles} transversaux } 2c^2 & 0 & 2 & 8 & 32
  \end{array}
\]
\note{« arêtes » est écrit au-dessus de « triangles » biffé}
\[
  2c^2=2\,2^{2\nu}=2^{2\nu+1}\qquad(\nu=-\infty,0,1,2)
\]
\[
  \begin{array}{l|c|c|c|c}
  \text{NB des sommets antiliés}
   & 0 & 4=2^2 & 8=2^3 & 16=2^4\\
  \text{à un sommet } 4c=2^{\nu+2} & & =2^e & =2^e & =2^{e-1}\\\hline
  \text{Ordre de } \mathrm{Aut}(S,A) : N_\omega & 6=2.3 & 72=2^3.3^2 & \cdot & 51840\\
   & & & & =2^7.3^4.5
  \end{array}
\]
\note{la case $c=2$ de l'ordre de $\mathrm{Aut}(S,A)$ est surchargée : on y lit un $360$ (?) recouvert et une factorisation en $2^{\cdot}.3^{\cdot}.5$ illisible}
\[
  \begin{array}{l|c|c|c|c}
  \text{Nb de } \tau_{-1}\ \text{(triangles) plongés}
   & 1 & 6=2.3 & 15=3.5 & 45=3^2 5\\
  \quad (c+1)(2c+1)=1+3c+2c^2 & & & & \\
  \text{Nb de } \tau_0\ \text{(bitriangles plongés)}
   & 0 & 1 & 10 & 120\\
  \quad \tfrac16\,c(c+1)(2c+1) & & & & \\
  \text{Nb de } \tau_1\ \text{(tritriangles) plongés}
   & 0 & 0 & 1 & 36=2^2 3^2\\
  \quad \tfrac1{30}\,c(c+1)(2c+1)(c-1) & & & &
  \end{array}
\]
\note{« tritriangles » est une lecture incertaine. Dans la ligne $\tau_0$, un « si $c=2^\nu$, $\nu\geqslant1,2$ » est biffé, et la valeur $120$ recouvre une écriture biffée ; le facteur de la ligne $\tau_{-1}$ est surchargé. Les formules $\frac16 c(c+1)(2c+1)$ et $\frac1{30}c(c+1)(2c+1)(c-1)$ sont telles quelles ; pour $c=4$ la seconde donne $18$, non le $36$ inscrit, que donne la formule de la p.~119}

comme \uncertain{les} figures précédentes, mais avec des \uncertain{$c\geqslant1$
épinglés} :
\[
  \begin{array}{c|c|c|c}
  6 & 36=2^2 3^2 & 90=2.3^2.5 & 270=2.3^3.5\\\hline
  0 & 72=2^3.3^2 & 360=2^3.3^2.5 & 2160=2^4.3^3.5\\\hline
  0 & 0 & 720=2^4.3^2 & 2^6.3^4.5\\\hline
  0 & 0 & - & 2^7.3^4.5
  \end{array}
\]
\note{dans ce dernier tableau, plusieurs cases sont surchargées ($36$ récrit sur une autre valeur, $2160$ ajouté au-dessus d'un nombre biffé) ; « $720=2^4.3^2$ » est tel quel. Les valeurs de la deuxième ligne diffèrent de celles du tableau des $c'_{0,\gamma}$ de la p.~118}
\note{en bas à droite : « $c=2^\nu$ »}

\end{document}
